\chapter{Системы счисления и коды}
\label{ch:codes}

В предыдущей главе мы записывали числа в двоичной системе. Для машины это
естественно, а человеку неудобно: число 200 превращается в
\code{11001000}. К тому же в цифровой технике кодируют не только числа, но и
буквы, цифры на индикаторе, положение вала, а ещё код может защищать данные от
ошибок. В этой главе мы разберём самые употребительные коды и
научимся выбирать подходящий.

\section{Шестнадцатеричная и восьмеричная системы}

Четыре двоичных разряда ровно соответствуют одной шестнадцатеричной цифре.
Поэтому длинные двоичные числа записывают группами по четыре бита, а каждую
группу заменяют одним символом: от \code{0} до \code{9}, затем \code{A}--\code{F}.

\begin{table}[H]
\centering
\caption{Соответствие двоичных групп и шестнадцатеричных цифр}
\small
\begin{tabular}{*{8}{c}}
\toprule
\textbf{2} & \textbf{16} & \textbf{2} & \textbf{16} & \textbf{2} & \textbf{16} & \textbf{2} & \textbf{16} \\
\midrule
0000 & 0 & 0100 & 4 & 1000 & 8 & 1100 & C \\
0001 & 1 & 0101 & 5 & 1001 & 9 & 1101 & D \\
0010 & 2 & 0110 & 6 & 1010 & A & 1110 & E \\
0011 & 3 & 0111 & 7 & 1011 & B & 1111 & F \\
\bottomrule
\end{tabular}
\end{table}

Например, $11001000_2 = \text{C8}_{16} = 12\cdot 16 + 8 = 200$. В Verilog
такие числа пишут с префиксом разрядности: \code{8'hC8}, \code{8'b11001000},
\code{8'd200}. Восьмеричная система, в которой одна цифра соответствует трём
битам, встречается реже: она осталась в правах доступа к файлам Unix
(\code{chmod 755}).

\begin{infobox}[Как быстро перевести]
Из двоичной в шестнадцатеричную: разбить число справа налево на четвёрки,
добавив слева нули, и заменить каждую четвёрку цифрой. Обратно --- развернуть
каждую цифру в четыре бита. Умножение и деление на степень двойки --- это
просто сдвиг: $x \cdot 2^k$ дописывает справа $k$ нулей.
\end{infobox}

\section{Двоично-десятичный код}

Человек привык к десятичной системе, поэтому устройства, которые показывают
числа (часы, калькуляторы, приборы), часто хранят не двоичное число, а
\emph{каждую десятичную цифру отдельно} в виде четырёх бит. Такой код
называется \term{двоично-десятичным} (BCD, Binary-Coded Decimal) в весах
8--4--2--1.

\begin{table}[H]
\centering
\caption{Число 2026 в разных видах}
\small
\begin{tabularx}{\textwidth}{l X}
\toprule
\textbf{Вид} & \textbf{Запись} \\
\midrule
Десятичное & 2\,0\,2\,6 \\
Двоичное & \code{0111 1110 1010} (12 бит) \\
Шестнадцатеричное & \code{7EA} \\
BCD & \code{0010 0000 0010 0110} (16 бит, по цифре в тетраде) \\
\bottomrule
\end{tabularx}
\end{table}

Достоинства BCD: вывод на индикатор сводится к простому дешифратору для одной
цифры (глава~\ref{ch:coders}), а округление десятичных дробей получается
точным. Недостатки: код расточителен (комбинации \code{1010}--\code{1111} не
используются), а арифметика сложнее. Если сложить две цифры BCD обычным
сумматором и получить результат больше девяти, его нужно исправить:
прибавить \code{0110} (шесть), чтобы пропустить шесть лишних комбинаций и
вызвать перенос в следующую цифру.

\begin{figure}[H]
\centering
\begin{tikzpicture}[x=1.0cm,y=1.0cm, font=\sffamily\small]
  \node[anchor=east] at (0,1) {$0111$};
  \node[anchor=east] at (0,0.4) {$+\ 0101$};
  \draw (-1.6,0.1) -- (0,0.1);
  \node[anchor=east] at (0,-0.3) {$1100$};
  \node[anchor=west, text=FpgaGray] at (0.4,-0.3) {$12>9$ --- не цифра BCD};
  \node[anchor=east] at (0,-0.9) {$+\ 0110$};
  \draw (-1.6,-1.2) -- (0,-1.2);
  \node[anchor=east] at (0,-1.6) {$1\,0010$};
  \node[anchor=west, text=FpgaGreen!60!black] at (0.4,-1.6) {перенос 1 и цифра 2: это $12$ в BCD};
  \node[anchor=east] at (0,2) {\textbf{7 + 5}};
\end{tikzpicture}
\caption{Сложение $7+5$ в BCD: после обычного сложения прибавляем поправку 6}
\end{figure}

\section{Коды символов}

Буквы тоже кодируют числами. Самый известный код --- \term{ASCII}: 7 бит на
символ, 128 символов. Латинские буквы идут подряд (\code{A} = 65 = \code{0x41},
\code{B} = 66 и далее), цифры тоже (\code{0} = 48 = \code{0x30}), поэтому
код цифры от кода символа отличается на $48$. Первые 32 кода --- управляющие
(перевод строки \code{0x0A}, возврат каретки \code{0x0D}). В Verilog символ в
апострофах превращается в его ASCII-код: \code{"A"} это то же самое, что
\code{8'h41}. Русские буквы и прочие символы в современных системах кодирует
UTF-8, где один символ занимает от одного до четырёх байт.

Код Грея мы уже встречали в главе~\ref{ch:coders}: соседние числа в нём
отличаются одним битом, поэтому датчик положения, который считывает код с
диска, не ошибётся больше чем на одну позицию. Он же пригодится нам в
главе~\ref{ch:clocks}, когда нужно передать указатель между двумя тактовыми
доменами.

\section{Обнаружение ошибок: бит чётности}

Любой канал связи и любая память иногда ошибаются: на проводе наводка, в
ячейке памяти пролетела космическая частица. Самый простой способ заметить
ошибку --- добавить к данным один бит так, чтобы число единиц во всём слове
стало чётным. Это \term{бит чётности}.

\begin{table}[H]
\centering
\caption{Бит чётности для 4-битных данных}
\small
\begin{tabular}{ccc}
\toprule
\textbf{Данные} & \textbf{Единиц} & \textbf{Бит чётности} \\
\midrule
0000 & 0 & 0 \\
0111 & 3 & 1 \\
1011 & 3 & 1 \\
1111 & 4 & 0 \\
\bottomrule
\end{tabular}
\end{table}

Бит чётности равен исключающему ИЛИ всех битов данных: $p = d_0 \oplus d_1
\oplus d_2 \oplus d_3$. Принимающая сторона считает ту же сумму по модулю два
уже вместе с битом чётности; если получилась единица, одна из позиций (или
нечётное их число) повредилась. Метод обнаруживает одну ошибку, но не знает,
где она, и не замечает двух ошибок сразу.

\section{Исправление ошибок: код Хэмминга}

Можно ли не только заметить, но и \emph{исправить} ошибку? Да, если
контрольных битов несколько и каждый проверяет свою часть слова. Классический
пример --- \term{код Хэмминга (7,4)}: четыре бита данных превращаются в
семь, и любая одиночная ошибка исправляется.

Позиции битов нумеруют от 1 до 7. На позициях, равных степеням двойки (1, 2,
4), стоят контрольные биты $p_1, p_2, p_3$, а на остальных --- данные
$d_1 \ldots d_4$. Каждый контрольный бит проверяет те позиции, в двоичной
записи номера которых есть соответствующая единица.

\begin{table}[H]
\centering
\caption{Позиции кода Хэмминга (7,4) и области проверки}
\small
\begin{tabular}{l*{7}{c}}
\toprule
\textbf{Позиция} & 1 & 2 & 3 & 4 & 5 & 6 & 7 \\
\textbf{Содержимое} & $p_1$ & $p_2$ & $d_1$ & $p_3$ & $d_2$ & $d_3$ & $d_4$ \\
\midrule
$p_1$ проверяет & \checkmark & & \checkmark & & \checkmark & & \checkmark \\
$p_2$ проверяет & & \checkmark & \checkmark & & & \checkmark & \checkmark \\
$p_3$ проверяет & & & & \checkmark & \checkmark & \checkmark & \checkmark \\
\bottomrule
\end{tabular}
\end{table}

Кодер вычисляет каждый $p$ как чётность своей группы:
$p_1 = d_1\oplus d_2\oplus d_4$, $p_2 = d_1\oplus d_3\oplus d_4$,
$p_3 = d_2\oplus d_3\oplus d_4$. Декодер по принятому слову пересчитывает три
проверки $s_1, s_2, s_3$. Если ошибки нет, все они равны нулю. Если ошибка
в позиции $k$, то набор $(s_3 s_2 s_1)$ как раз является двоичной записью
числа $k$. Эта тройка называется \term{синдромом}: она прямо указывает на
позицию ошибки, и остаётся только перевернуть этот бит.

\begin{ideabox}
Например, данные $d_1d_2d_3d_4 = 1011$. Кодер даёт $p_1 = 1\oplus 0\oplus 1 = 0$,
$p_2 = 1\oplus 1\oplus 1 = 1$, $p_3 = 0\oplus 1\oplus 1 = 0$, и в канал уходит слово
\code{0110011}. Если шестой бит перевернётся, декодер получит $s_1=0$, $s_2=1$,
$s_3=1$, то есть синдром $110_2=6$. Бит на позиции 6 переворачивается обратно, и
данные восстановлены.
\end{ideabox}

Опишем кодер, канал с ошибкой и декодер на Verilog. Получится схема, целиком
состоящая из исключающих ИЛИ.

\vfile{\codepath/00_theory/hamming74.v}{code/00\_theory/hamming74.v}

\shot{n-hamming-rtl}{Код Хэмминга в Yosis: деревья исключающего ИЛИ вычисляют контрольные биты и синдром}{fig:n-hamming}

Если в вашем проекте есть проверяющий модуль, его удобно тоже написать на
Verilog (файл \code{tb\_hamming74.v} лежит рядом): он перебирает все 16
значений данных и все семь одиночных ошибок и убеждается, что данные
восстанавливаются каждый раз.

Такой код применяют в памяти серверов (ECC), в линиях связи, в
QR-кодах (там используют более мощный код Рида---Соломона). Чем больше
контрольных битов, тем больше ошибок можно исправить, но тем больше
избыточность. Для 64 бит данных хватает 8 контрольных (7 по схеме Хэмминга и один общий бит чётности для обнаружения двойных ошибок).

\section{Дробные числа}

Целые числа хранить просто. А как хранить $3{,}14$? Есть два способа.

\term{Фиксированная запятая.} Условно считаем, что запятая стоит
между фиксированными разрядами. В формате Q4.4 байт делится пополам: старшие
4 бита --- целая часть, младшие 4 --- дробная с весами $\tfrac12, \tfrac14,
\tfrac18, \tfrac1{16}$. Тогда \code{0011\,0100} это $3 + \tfrac14 = 3{,}25$. Все
операции сложения и вычитания выполняются обычным сумматором, а
умножение --- обычным умножителем, после которого результат сдвигается на
число дробных разрядов. Для сигнальной обработки это основной способ: он
экономит ресурсы ПЛИС.

\term{Плавающая запятая.} Число хранится как $(-1)^s \cdot m \cdot
2^{e}$: знак, мантисса и порядок. Стандарт IEEE~754 для 32-битного формата
\code{float} отводит 1 бит на знак, 8 на порядок и 23 на мантиссу. Диапазон
огромен (порядка $10^{\pm 38}$), но точность ограничена семью десятичными
знаками, а число $0{,}1$ в двоичной системе вообще не представимо точно (оно
бесконечная периодическая дробь), как $\tfrac13$ в десятичной. Блок
сложения чисел с плавающей запятой гораздо сложнее и больше целочисленного,
поэтому в небольших ПЛИС его стараются не использовать.

\begin{table}[H]
\centering
\caption{Выбор представления числа}
\small
\begin{tabularx}{\textwidth}{l X X}
\toprule
\textbf{Код} & \textbf{Где применяют} & \textbf{Особенность} \\
\midrule
Двоичный & Счётчики, адреса, арифметика & Самый компактный \\
Дополнительный & Числа со знаком & Вычитание сводится к сложению \\
BCD & Индикаторы, часы, финансы & Просто выводить, точная десятичная дробь \\
Код Грея & Датчики угла, FIFO между доменами & Соседние значения отличаются одним битом \\
Бит чётности & Простые каналы & Находит одиночную ошибку \\
Код Хэмминга & Память, связь & Исправляет одиночную ошибку \\
Фиксированная запятая & Цифровая обработка сигналов & Быстро и дёшево \\
Плавающая запятая & Научные расчёты & Огромный диапазон, но дорого \\
\bottomrule
\end{tabularx}
\end{table}

\section{Вопросы и задачи}
\label{sec:tasks-codes}

Ответы собраны в приложении~\ref{app:answers}.

\begin{taskbox}[Задачи к главе]
\begin{enumerate}
  \item Переведите $\text{B7}_{16}$ и $\text{3C}_{16}$ в двоичную и десятичную системы.
  \item Запишите число 1998 в двоично-десятичном коде.
  \item Закодируйте данные $d_1d_2d_3d_4 = 0110$ кодом Хэмминга (7,4).
  \item Принято слово \code{1100111}. Найдите синдром, укажите повреждённую позицию и восстановите данные.
  \item Данные \code{1011001}. Чему равен бит чётности (дополняем до чётного числа единиц)?
  \item Чему равно число \code{0x2A} в формате с фиксированной запятой Q4.4?
\end{enumerate}
\end{taskbox}
